% ==========================================================
% REFERENCIAS
% Añade cada fuente nueva como un \item más, numerada por
% orden de primera aparición en el texto.
% ==========================================================

\section*{\centering REFERENCIAS}


\begin{enumerate}[leftmargin=1.8em,label=\lbrack\arabic*\rbrack,itemsep=0.6em]

\item F. Ccami-Bernal, P. Urday-Ramos, F. Zela-Coila, J. K. Cáceres-Ruiz y V. Cabrera-Caso, ``Trastornos musculoesqueléticos y prácticas ergonómicas en universitarios peruanos durante la pandemia de la COVID-19,'' \emph{Revista Cubana de Ortopedia y Traumatología}, vol. 37, no. 3, 2023. [En línea]. Disponible en: \url{https://revortopedia.sld.cu/index.php/revortopedia/article/view/608}

\item D. Fernández-Villacorta, A. N. Soriano-Moreno, T. Gálvez-Olórtegui, N. Agüero-Santivañez, D. R. Soriano-Moreno y V. A. Benites-Zapata, ``Síndrome visual informático en estudiantes universitarios de posgrado de una universidad privada de Lima, Perú,'' \emph{Archivos de la Sociedad Española de Oftalmología}, vol. 96, no. 10, pp. 515--520, 2021. doi: 10.1016/j.oftal.2020.12.003. [En línea]. Disponible en: \url{https://www.elsevier.es/es-revista-archivos-sociedad-espanola-oftalmologia-296-articulo-sindrome-visual-informatico-estudiantes-universitarios-S0365669121000058}

\item R. M. Meneses Castañeda, S. L. Ramos Rodríguez, C. D. C. Molfino Jaramillo, E. L. Sánchez Miraval, D. F. Stein Montoros y L. G. Chávez Rodríguez, ``Síndrome visual informático en estudiantes de medicina en educación virtual de una universidad peruana durante el 2021,'' \emph{Revista de la Facultad de Medicina Humana}, vol. 23, no. 1, pp. 25--32, 2023. doi: 10.25176/RFMH.v23i1.5525. [En línea]. Disponible en: \url{https://doi.org/10.25176/RFMH.v23i1.5525}

\item Zhunio Suin et al., ``Pausas activas y su importancia en el bienestar de estudiantes universitarios,'' 2024.

\item L. C. Gerena Pallares, L. J. Vargas Rodríguez, C. A. Niño Avendaño, G. C. Uyaban y Y. Ballesteros Virgen, ``Prevalencia del síndrome visual por computadora en los estudiantes de medicina de la ciudad de Tunja durante la pandemia,'' \emph{Revista Colombiana de Salud Ocupacional}, vol. 12, no. 1, e-7916, jul. 2022. doi: 10.18041/2322-634X/rcso.1.2022.7916. [En línea]. Disponible en: \url{https://doi.org/10.18041/2322-634X/rcso.1.2022.7916}

\item A. J. Almario Barrera, C. A. Orgulloso Bautista, A. del C. Padrón Merlano, J. A. Díaz Cetina y Y. Z. Castellanos Domínguez, ``Factores asociados a trastornos musculoesqueléticos entre universitarios del nororiente colombiano durante la pandemia por COVID-19,'' \emph{Universitas Odontológica}, vol. 42, 2023. doi: 10.11144/Javeriana.uo42.famd. [En línea]. Disponible en: \url{https://doi.org/10.11144/Javeriana.uo42.famd}

\item E. E. García-Salirrosas y R. A. Sánchez-Poma, ``Prevalencia de los trastornos musculoesqueléticos en docentes universitarios que realizan teletrabajo en tiempos de COVID-19,'' \emph{Anales de la Facultad de Medicina}, vol. 81, no. 3, pp. 301--307, 2020. doi: 10.15381/anales.v81i3.18841. [En línea]. Disponible en: \url{https://doi.org/10.15381/anales.v81i3.18841}

\item C. L. Ruiz Sanabria, E. Y. Niño Vega y S. P. Jurado, ``Correlación entre la sintomatología ocular asociada al síndrome visual informático e insuficiencia de convergencia,'' \emph{Ciencia y Tecnología para la Salud Visual y Ocular}, vol. 21, no. 1, 2023. doi: 10.19052/sv.vol21.iss1.6. [En línea]. Disponible en: \url{https://doi.org/10.19052/sv.vol21.iss1.6}

\item M. Sánchez, S. Rojas, J. F. Cuarán, K. A. García, D. E. Palacio, N. L. Valencia y M. O. Garzón, ``Factores asociados a afecciones oftalmológicas en un grupo de estudiantes universitarios de la ciudad de Medellín, prepandemia y pandemia, 2021,'' \emph{Medicina U.P.B.}, vol. 43, no. 2, pp. 26--34, 2024. doi: 10.18566/medupb.v43n2.a04. [En línea]. Disponible en: \url{https://doi.org/10.18566/medupb.v43n2.a04}

\item K. Quilumba-Gualoto, ``Factores de riesgo e intervenciones ergonómicas efectivas para el manejo del síndrome de visión de computadora,'' \emph{Ergonomía, Investigación y Desarrollo}, vol. 1, no. 3, pp. 10--21, 2019. [En línea]. Disponible en: \url{https://revistas.udec.cl/index.php/Ergonomia_Investigacion/article/download/1345/1974/2396}

\item R. Hernández-Sampieri, C. Fernández-Collado and P. Baptista-Lucio,
\textit{Metodología de la investigación}, 6th ed.
México D.F., México: McGraw-Hill, 2014.

\item Y. Torres y Y. Rodríguez, ``Surgimiento y evolución de la ergonomía como disciplina: reflexiones sobre la escuela de los factores humanos y la escuela de la ergonomía de la actividad,'' \emph{Revista Facultad Nacional de Salud Pública}, vol. 39, no. 2, e342868, 2021. doi: 10.17533/udea.rfnsp.e342868. [En línea]. Disponible en: \url{https://doi.org/10.17533/udea.rfnsp.e342868}

\item L. Chicata, D. E. Palomino-Pagador y F. M. Runzer-Colmenares, ``Tiempo de exposición a pantallas digitales y su asociación con el síndrome visual informático en estudiantes de medicina,'' \emph{Acta Médica Peruana}, vol. 41, no. 4, pp. 250--258, 2024. doi: 10.35663/amp.2024.414.3250. [En línea]. Disponible en: \url{https://doi.org/10.35663/amp.2024.414.3250}

\item A. K. Lema y E. W. Anbesu, ``Computer vision syndrome and its determinants: A systematic review and meta-analysis,'' \emph{SAGE Open Medicine}, vol. 10, 2022. doi: 10.1177/20503121221142402. [En línea]. Disponible en: \url{https://doi.org/10.1177/20503121221142402}

\item T. J. Armstrong et~al., ``A conceptual model for work-related neck and upper-limb musculoskeletal disorders,'' \emph{Scandinavian Journal of Work, Environment \& Health}, vol. 19, no. 2, pp. 73--84, 1993. doi: 10.5271/sjweh.1494. [En línea]. Disponible en: \url{https://doi.org/10.5271/sjweh.1494}

\item M. Rosenfield, ``Computer vision syndrome: a review of ocular causes and potential treatments,'' \emph{Ophthalmic and Physiological Optics}, vol. 31, no. 5, pp. 502--515, 2011. doi: 10.1111/j.1475-1313.2011.00834.x. [En línea]. Disponible en: \url{https://doi.org/10.1111/j.1475-1313.2011.00834.x}

\item C. Blehm, S. Vishnu, A. Khattak, S. Mitra y R. W. Yee, ``Computer vision syndrome: a review,'' \emph{Survey of Ophthalmology}, vol. 50, no. 3, pp. 253--262, 2005. doi: 10.1016/j.survophthal.2005.02.008. [En línea]. Disponible en: \url{https://doi.org/10.1016/j.survophthal.2005.02.008}

\item Congreso de la República de Colombia, ``Ley 1581 de 2012, por la cual se dictan disposiciones generales para la protección de datos personales,'' Bogotá, Colombia: Diario Oficial, 2012.

\item World Medical Association, ``WMA Declaration of Helsinki: ethical principles for medical research involving human participants.'' [En línea]. Disponible en: \url{https://www.wma.net/policies-post/wma-declaration-of-helsinki/}

\item Ministerio de Salud de Colombia, ``Resolución 8430 de 1993, por la cual se establecen las normas científicas, técnicas y administrativas para la investigación en salud,'' Bogotá, Colombia, 1993.

\item M. d. M. Seguí, J. Cabrero-García, A. Crespo, J. Verdú y E. Ronda, ``A reliable and valid questionnaire was developed to measure computer vision syndrome at the workplace,'' \emph{Journal of Clinical Epidemiology}, vol. 68, no. 6, pp. 662--673, 2015. doi: 10.1016/j.jclinepi.2015.01.015. [En línea]. Disponible en: \url{https://doi.org/10.1016/j.jclinepi.2015.01.015}

\item D. C. Silva Sánchez, D. F. Rojas-Gualdrón, N. E. Gómez Rúa, E. Ronda Pérez y M. d. M. Seguí Crespo, ``Adaptación cultural y validación del Computer Vision Syndrome Questionnaire en Colombia,'' \emph{Revista Cuidarte}, vol. 16, no. 2, e3963, 2025. doi: 10.15649/cuidarte.3963. [En línea]. Disponible en: \url{https://doi.org/10.15649/cuidarte.3963}

\item I. Kuorinka et~al., ``Standardised Nordic questionnaires for the analysis of musculoskeletal symptoms,'' \emph{Applied Ergonomics}, vol. 18, no. 3, pp. 233--237, 1987. doi: 10.1016/0003-6870(87)90010-X.

\item L. J. Cronbach, ``Coefficient alpha and the internal structure of tests,'' \emph{Psychometrika}, vol. 16, no. 3, pp. 297--334, 1951.

\end{enumerate}


